\documentclass[conference]{IEEEtran}
\IEEEoverridecommandlockouts
\usepackage{cite}
\usepackage{amsmath,amssymb,amsfonts}
\usepackage{algorithmic}
\usepackage{graphicx}
\usepackage{textcomp}
\usepackage{xcolor}
\usepackage{booktabs}
\usepackage{url}
\usepackage{balance}

\def\BibTeX{{\rm B\kern-.05em{\sc i\kern-.025em b}\kern-.08em
    T\kern-.1667em\lower.7ex\hbox{E}\kern-.125emX}}

\begin{document}

\title{ResiliPay: A Fault-Tolerant, Dual-Channel Cryptographic Payment Protocol for Intermittently Connected Economies\thanks{Archived on Zenodo (CERN/OpenAIRE): \url{https://doi.org/10.5281/zenodo.23451943}}}

\author{
\IEEEauthorblockN{Brayan Osinaka}
\IEEEauthorblockA{\textit{National Institute of Information Technology} \\
Rivers State, Nigeria \\
brayanjunior101@gmail.com}
}

\maketitle

\begin{abstract}
Digital payment adoption across emerging economies has surged, yet contemporary architectures remain fundamentally brittle due to their strict reliance on continuous, low-latency telecommunications and synchronous centralized clearing switches. In countries such as Nigeria, recurring electrical grid collapses and cellular backhaul disruptions cause transaction failure rates between 10\% and 15\% during peak retail hours, imposing over $\mathbb{N}$100 billion in annual deadweight losses across distributed agency banking networks. Existing offline proposals suffer from a fundamental trilemma: they either permit catastrophic offline double-spending, introduce delayed-settlement insolvency risks (unbacked overdrafts), or demand specialized, expensive smartcard hardware.

This paper presents \textbf{ResiliPay}, an asynchronous, offline-first cryptographic payment protocol designed to sustain retail commerce during total wide-area network and telecommunication partitions. ResiliPay combines client-side asymmetric cryptography (ECDSA P-256) with hardware-enforced monotonic counters and a Pre-Allocated Offline Escrow (POE) framework hosted inside mobile Trusted Execution Environments (TEEs) to mathematically eliminate both offline double-spending and delayed-settlement insolvency without requiring distributed consensus. To bridge physical transport gaps under infrastructure failures, ResiliPay introduces an adaptive dual-channel serialization scheme: high-density optical Quick Response (QR) codes compacted by 72.4\% (125 bytes, QR Version 6/8) for zero-connectivity proximity transactions, and an ultra-compact bit-packed binary encoding (86 bytes) that fits within a single GSM 7-bit SMS segment ($<$160 characters) for wide-area cellular fallback. 

Furthermore, ResiliPay implements an order-independent reconciliation state machine utilizing sliding replay-window nonce sets, guaranteeing that out-of-order merchant batch uploads after regional grid restoration achieve 100\% settlement precision without false replay rejections. Under a simulated post-blackout reconnection burst of 1,000 out-of-order merchant transactions with 50 injected adversarial replay attacks, ResiliPay achieves 100\% settlement precision and 100\% adversarial rejection in 176\,ms ($>$5,600 settled tx/sec). We outline a pragmatic three-tier hardware security model spanning dedicated Secure Elements, ARM TrustZone with Replay Protected Memory Blocks (RPMB), and bank-underwritten micro-floats, providing an actionable blueprint for sovereign financial resilience in infrastructure-challenged economies.
\end{abstract}

\begin{IEEEkeywords}
Offline Payments, Financial Inclusion, ICT4D, Hardware Security Modules, ARM TrustZone, ECDSA, Dual-Channel Communication, Intermittent Connectivity, Delay-Tolerant Systems.
\end{IEEEkeywords}

\section{Introduction}
The transition from physical cash to sovereign electronic transaction systems represents one of the most transformative socio-economic developments across the Global South~\cite{aker2010mobile, jack2014risk}. In Sub-Saharan Africa, mobile money platforms and instant payment switches have brought tens of millions of formerly unbanked citizens into formal financial systems~\cite{gsma2023state}. In Nigeria, the Nigeria Inter-Bank Settlement System Instant Payment (NIP) switch processed electronic transfers valued at over \$600 billion in 2023, driven largely by a distributed footprint of over 1.5 million agency banking retail agents~\cite{nibss2023annual, efina2023access}.

Despite rapid adoption, contemporary retail digital payment architectures suffer from a critical architectural vulnerability: \textbf{synchronous network dependency}. Existing card payment rails (EMV), point-of-sale (POS) terminals, and mobile banking applications demand an active, low-latency bi-directional round-trip connection between consumer smartphones, merchant terminals, acquirer gateways, centralized switches, and issuing bank core ledgers~\cite{bis2023polaris, chu2022offline}. 

In emerging economies, this assumption of uninterrupted telecommunications fails systematically:
\begin{enumerate}
    \item \textbf{Electrical Grid Fragility:} Periodic collapses of national electrical grids cause cellular base transceiver stations (BTS) to exhaust emergency diesel generator and battery reserves, triggering widespread cellular blackouts lasting from hours to days.
    \item \textbf{Switch Concurrency Cascades:} During peak market hours, sudden surges in transaction volumes overload banking core databases, causing timeout cascades and peak failure rates between 10\% and 15\%~\cite{cbn2023stability}.
    \item \textbf{Economic Exclusion in Agency Banking:} Agency banking retail operators in semi-urban and rural markets are routinely forced to turn away customers or revert to scarce physical cash when POS terminals display ``Switch Unavailable'' or ``Network Timeout.''
\end{enumerate}

\subsection{The Offline Payment Trilemma}
Developing an offline digital payment system involves navigating a well-known architectural trilemma:
\begin{enumerate}
    \item \textbf{Double-Spending Vulnerability:} Purely software-based systems cannot prevent an offline user from spending the same balance across multiple isolated merchants without synchronous consensus.
    \item \textbf{Delayed-Settlement Insolvency (The Empty Account Dilemma):} While digital signatures prove voucher authorization, an offline merchant cannot verify whether the consumer's bank account actually contains sufficient funds. If a malicious consumer with $\mathbb{N}$500 in their account issues multiple $\mathbb{N}$50,000 offline vouchers to different merchants, the clearing switch will discover insolvency upon reconnection, imposing catastrophic credit default losses on merchants.
    \item \textbf{Hardware Cost and Friction:} Solutions relying on proprietary smartcards (e.g., Mondex) or specialized tamper-resistant POS terminals impose prohibitive capital expenditure costs on low-income retail agents.
\end{enumerate}

\subsection{Key Contributions}
To resolve this trilemma, this paper introduces \textbf{ResiliPay}, a fault-tolerant, decentralized offline payment protocol:
\begin{itemize}
    \item \textbf{Hardware-Enforced Double-Spend \& Insolvency Pre-emption:} We couple Elliptic Curve Digital Signature Algorithm (ECDSA P-256) signatures with hardware-enforced monotonic counters and a Pre-Allocated Offline Escrow (POE) mechanism situated within commodity mobile Trusted Execution Environments (TEEs), eliminating both double-spending and insolvency defaults without requiring active network connectivity.
    \item \textbf{Dual-Channel Physical Transmission:} We design an adaptive serialization format supporting optical high-density QR transmission (compacted by 72.4\% over naive JSON representations to 125 bytes, QR Version 6/8) and a compressed 86-byte single-segment GSM 7-bit SMS fallback ($<$160 characters), avoiding multi-part SMS fees.
    \item \textbf{Order-Independent Reconciliation Engine:} We formalize an asynchronous clearing state machine utilizing sliding replay-window nonce sets. We prove and demonstrate via simulation that 1,000 transactions submitted by 50 merchants in completely randomized, out-of-order batches achieve 100\% settlement precision and 100\% adversarial rejection in 176\,ms ($>$5,600 tx/sec throughput).
    \item \textbf{Pragmatic Three-Tier Hardware Security Model:} We define a tiered deployment architecture spanning dedicated Secure Elements, ARM TrustZone with Replay Protected Memory Blocks (RPMB) on commodity MediaTek/Unisoc chipsets, and bank-underwritten micro-floats, establishing realistic feasibility for low-cost smartphones in Sub-Saharan Africa.
\end{itemize}

---

\section{System Architecture \& Threat Model}

\subsection{System Actors}
The ResiliPay architecture comprises three primary entities:
\begin{enumerate}
    \item \textbf{Payer (Consumer Device, $\mathcal{C}$):} An entry-level or mid-range smartphone running a secure mobile banking client equipped with an on-device hardware TEE.
    \item \textbf{Payee (Merchant / Agency Banking Terminal, $\mathcal{M}$):} A smartphone or Android POS terminal running an offline verification engine capable of scanning optical QR matrices or receiving SMS records.
    \item \textbf{Settlement Core (Central Bank / Clearinghouse, $\mathcal{S}$):} The centralized switch (e.g., NIBSS or commercial bank core) that ingests queued transaction batches upon network restoration and updates master account balances.
\end{enumerate}

\subsection{Threat Model and Security Objectives}
We consider an adversarial environment where wide-area telecommunications are totally partitioned, meaning neither $\mathcal{C}$ nor $\mathcal{M}$ can communicate with $\mathcal{S}$ during transaction execution.

\subsubsection{Adversarial Capabilities}
\begin{itemize}
    \item \textbf{Network Adversary:} An attacker capable of intercepting, dropping, delaying, or tampering with wireless transmissions across optical, cellular, or SMS channels.
    \item \textbf{Malicious Payer ($\mathcal{C}^*$):} A dishonest consumer who may attempt to:
    \begin{enumerate}
        \item Alter transaction amounts or merchant identifiers after signature generation.
        \item Replay a spent voucher to another offline merchant (\textit{Replay Attack}).
        \item Fork their state to generate two conflicting vouchers with identical monotonic nonces (\textit{Offline Double-Spending Attack}).
        \item Issue offline vouchers that cumulatively exceed their actual liquid bank balance (\textit{Insolvency / Zero-Balance Default Attack}).
        \item Roll back the mobile operating system or re-flash firmware to reset account balances.
    \end{enumerate}
    \item \textbf{Untrusted Merchant ($\mathcal{M}^*$):} A dishonest merchant attempting to duplicate a valid customer voucher to claim multiple credits during batch settlement.
\end{itemize}

\subsubsection{Security Objectives}
\begin{itemize}
    \item \textbf{Authenticity \& Non-Repudiation:} A merchant must be guaranteed that the voucher was authorized by the authenticated account holder and cannot be modified in transit.
    \item \textbf{Zero-Default Insolvency Guarantee:} Merchants must be mathematically or contractually protected against customer account insolvency upon delayed settlement.
    \item \textbf{Order-Independent Eventual Consistency:} Legitimate merchant vouchers must settle successfully regardless of the temporal order in which merchant queues reconnect to $\mathcal{S}$.
\end{itemize}

---

\section{The ResiliPay Protocol}

\subsection{Cryptographic Primitives and Hardware Enclave Binding}
ResiliPay utilizes NIST P-256 (secp256r1) elliptic curve cryptography paired with SHA-256 hashing~\cite{nist2013fips186}.

\subsubsection{Onboarding \& Key Attestation}
During initial online enrollment with the financial institution:
\begin{equation}
\mathcal{G}_{\text{TEE}}() \rightarrow (sk_{\mathcal{C}}, pk_{\mathcal{C}}, \Pi_{\text{attest}})
\end{equation}
The private key $sk_{\mathcal{C}}$ is generated inside the mobile device's hardware enclave (ARM TrustZone or dedicated Secure Element)~\cite{android2024keystore, apple2023secureenclave}. The key is bound to non-exportable hardware flags:
\begin{equation}
\begin{split}
\text{Flags} = \{ & \text{PURPOSE\_SIGN}, \\
                  & \text{NON\_EXPORTABLE}, \\
                  & \text{HARDWARE\_COUNTER} \}
\end{split}
\end{equation}
The hardware issues an asymmetric Key Attestation certificate $\Pi_{\text{attest}}$ signed by the silicon manufacturer root of trust, certifying to the issuing bank that $sk_{\mathcal{C}}$ is hardware-isolated and incapable of state extraction or counter reset.

\subsection{Pre-Allocated Offline Escrow (POE) Framework}
To mathematically eliminate delayed-settlement insolvency, ResiliPay enforces a \textbf{Pre-Allocated Offline Escrow (POE)}:
\begin{enumerate}
    \item While online, the consumer reserves an offline spending allowance $B_{\text{offline}}$ (e.g., $\mathbb{N}$20,000) from their primary bank balance into an escrow vault held at $\mathcal{S}$.
    \item The bank issues a cryptographically signed state token to the device TEE confirming the allocated escrow balance and an expiration epoch $T_{\text{exp}}$ (e.g., 7 days).
    \item The device TEE maintains an internal secure state tuple:
    \begin{equation}
    \mathcal{S}_{\text{TEE}} = \langle c, B_{\text{offline}}, T_{\text{exp}} \rangle
    \end{equation}
    where $c$ is the hardware monotonic counter.
    \item For micro-transactions below a bank-defined threshold (e.g., $\mathbb{N}$2,000), the issuing bank may alternatively issue an underwritten micro-float tied to the user's verified identity (BVN/NIN), contractually absorbing credit risk.
\end{enumerate}

\smallskip
\noindent\textbf{Epoch-Bound Leases \& Recovery from Lost Devices:}
If a consumer's device is physically destroyed, stolen, or drowned in water while unspent escrow remains inside $\mathcal{S}_{\text{TEE}}$, the consumer reports the incident to the issuing bank. The clearinghouse $\mathcal{S}$ places the remaining $B_{\text{offline}}$ into a quarantine window until $T_{\text{exp}} + \Delta_{\text{sync}}$ (e.g., 72 hours following regional network restoration). This cooling-off duration allows honest merchants to flush their offline queues $\mathcal{Q}_{\mathcal{M}}$. Once the epoch window expires, any unsettled residual escrow automatically sweeps back into the user's primary liquid account, pre-empting both fraudulent duplicate reporting and permanent capital loss.

\subsection{Offline Voucher Creation (Payer)}
When wide-area connectivity is absent, the payer initiates a transaction of amount $a$ to merchant $\mathcal{M}$ (identifier $ID_{\mathcal{M}}$):
\begin{enumerate}
    \item \textbf{Enclave Balance Verification:} The TEE verifies that $a \le B_{\text{offline}}$ and that current time $t \le T_{\text{exp}}$. If false, the transaction is rejected locally.
    \item \textbf{State Update:} The TEE atomically decrements $B_{\text{offline}} \leftarrow B_{\text{offline}} - a$ and increments its internal monotonic counter $c \leftarrow c + 1$.
    \item \textbf{Payload Construction:} The application builds the transaction message:
    \begin{equation}
    m = \langle ID_{\mathcal{C}}, ID_{\mathcal{M}}, a, c, t \rangle
    \end{equation}
    where $t$ represents the local Unix epoch timestamp.
    \item \textbf{Signing:} The TEE signs the SHA-256 digest using $sk_{\mathcal{C}}$:
    \begin{equation}
    \sigma = \text{Sign}_{sk_{\mathcal{C}}}(\mathcal{H}(m))
    \end{equation}
    \item The complete voucher package is assembled:
    \begin{equation}
    \mathcal{V} = \langle m, \sigma, pk_{\mathcal{C}} \rangle
    \end{equation}
\end{enumerate}

\subsection{Dual-Channel Physical Transmission}

\subsubsection{Primary Channel: High-Density Optical QR Code}
For in-person retail points, $\mathcal{C}$ renders $\mathcal{V}$ as a 2D Quick Response matrix. Naive JSON-in-Base64 serialization requires $\sim$604 bytes, forcing QR Version 14+ matrices (73$\times$73 modules) that suffer from severe scanning failure rates on cracked or low-resolution cameras.

ResiliPay introduces a structured binary layout (\textbf{ResiliPay-Compact}):
\begin{itemize}
    \item Timestamp $t$: 8 bytes (64-bit unsigned integer)
    \item Amount $a$: 4 bytes (32-bit unsigned integer, Kobo resolution)
    \item Monotonic Counter $c$: 4 bytes (32-bit unsigned integer)
    \item Payer $ID_{\mathcal{C}}$: 6 bytes (Binary-Coded Decimal for E.164 phone)
    \item Payee $ID_{\mathcal{M}}$: 6 bytes (Binary-Coded Decimal)
    \item Compressed Public Key: 33 bytes ($0x02/0x03 + \text{coord}_X$)
    \item ECDSA Signature $\sigma$: 64 bytes (IEEE P1363 $r \parallel s$)
\end{itemize}
\textbf{Total Binary Size: 125 Bytes.} 
In standard ISO/IEC 18004 QR encoding, 125 raw 8-bit bytes map into \textbf{QR Version 6} ($41\times41$ modules at Level L, 7\% error correction, max 134 bytes) or \textbf{QR Version 8} ($49\times49$ modules at Level M, 15\% error correction, max 152 bytes). Alternatively, using Base45 encoding (as utilized in international digital certificates), 125 bytes encode into 188 alphanumeric characters, fitting cleanly within \textbf{QR Version 7} ($45\times45$ modules at Level M). Compared to naive JSON Version 14 ($73\times73$ modules), this reduces matrix surface area by over \textbf{62\% to 68\%}, dramatically accelerating optical acquisition on low-end cameras under harsh ambient glare.

\subsubsection{Secondary Channel: Single-Segment GSM 7-bit SMS Fallback}
When optical scanning is unavailable (e.g., feature phones or broken lenses), ResiliPay provides an SMS signaling fallback:
\begin{equation}
m_{\text{SMS}} = \langle c \text{ (4B)}, a \text{ (4B)}, ID_{\mathcal{M}} \text{ (6B)}, t \text{ (4B)}, ID_{pk} \text{ (4B)}, \sigma \text{ (64B)} \rangle
\end{equation}
The total raw binary size of $m_{\text{SMS}}$ is \textbf{86 bytes}. In standard Base64 encoding, this constitutes 116 ASCII characters, fitting comfortably below the \textbf{160-character ceiling of a single GSM 7-bit SMS segment}.
Crucially, $m_{\text{SMS}}$ explicitly includes the 4-byte timestamp $t$ (ensuring mathematical reconstruction of $\mathcal{H}(m)$ during verification) and a 4-byte truncated public key hash prefix $ID_{pk}$ so merchant POS terminals can index locally cached customer public keys.
\begin{itemize}
    \item \textbf{Direct Cellular Relay:} The payer transmits $m_{\text{SMS}}$ to the bank gateway number; when cell towers recover, the switch clears the transaction and sends a USSD/SMS receipt to the merchant.
    \item \textbf{Proximity Relay:} When sent device-to-device, $ID_{pk}$ enables the merchant POS to look up locally cached customer keys from prior interactions.
\end{itemize}

\subsection{Offline Verification (Payee / Merchant)}
Upon receiving voucher $\mathcal{V}$ via optical scan or SMS:
\begin{enumerate}
    \item The merchant POS extracts $m$, $\sigma$, and $pk_{\mathcal{C}}$.
    \item \textbf{Cryptographic Verification:} The merchant executes:
    \begin{equation}
    \text{Verify}_{pk_{\mathcal{C}}}(\mathcal{H}(m), \sigma) \stackrel{?}{=} 1
    \end{equation}
    If mathematically invalid, the voucher is discarded.
    \item \textbf{Local Nonce Validation:} $\mathcal{M}$ checks its local offline ledger $\mathcal{L}_{\mathcal{M}}$ to confirm that counter $c$ has not been previously received from sender $ID_{\mathcal{C}}$.
    \item \textbf{Voucher Queuing:} If valid, $\mathcal{M}$ stores $\mathcal{V}$ in its offline sync queue $\mathcal{Q}_{\mathcal{M}}$ and issues goods immediately.
\end{enumerate}

\subsection{Order-Independent Settlement \& Reconciliation Engine}
Existing sequential ceiling checks ($c_i > \text{LedgerCounter}$) fail when multiple merchants reconnect out of chronological sequence. ResiliPay resolves this via an \textbf{Order-Independent Nonce Set Engine}:

When merchant $\mathcal{M}$ reconnects to wide-area connectivity:
\begin{enumerate}
    \item $\mathcal{M}$ submits queue $\mathcal{Q}_{\mathcal{M}} = \{\mathcal{V}_1, \mathcal{V}_2, \dots, \mathcal{V}_k\}$ to central switch $\mathcal{S}$.
    \item For each voucher $\mathcal{V}_i$, $\mathcal{S}$ validates:
    \begin{equation}
    \text{Verify}_{pk_{\mathcal{C}}}(\mathcal{H}(m_i), \sigma_i) = 1 \quad \land \quad c_i \notin \text{SettledCounters}[ID_{\mathcal{C}}]
    \end{equation}
    where $\text{SettledCounters}[ID_{\mathcal{C}}]$ is a persistent set of confirmed nonces within the active epoch.
    \item If valid, $\mathcal{S}$ commits:
    \begin{equation}
    \text{SettledCounters}[ID_{\mathcal{C}}] \leftarrow \text{SettledCounters}[ID_{\mathcal{C}}] \cup \{c_i\}
    \end{equation}
    and atomically credits amount $a_i$ from the payer's escrow balance to merchant $\mathcal{M}$.
\end{enumerate}
To prevent unbounded memory growth at the central switch, $\mathcal{S}$ maintains a contiguous low-water mark $c_{\min}$ (below which all counters are settled) and a compact 256-bit sliding bitmap for out-of-order nonces above $c_{\min}$ (analogous to IPsec anti-replay windows, RFC 4303). The bitmap is flushed upon POE epoch expiration, guaranteeing strict $O(1)$ RAM overhead per user at national scale.

---

\section{Security Analysis}

\subsection{Mathematical Impossibility of Offline Double-Spending}
\textbf{Theorem 1 (Counter Monotonicity):} Let $c_k$ denote the $k$-th hardware monotonic counter state. For all $j > k$, $c_j > c_k$. An attacker possessing root privileges over the host mobile operating system cannot reset $c_k$ without triggering a silicon-level cryptographic zeroization of the private key $sk_{\mathcal{C}}$.

\textbf{Theorem 2 (Uniqueness of Signed State):} Because the hardware counter increments on every signature invocation and is bound to $m$, an attacker cannot generate two distinct valid vouchers $\mathcal{V}_1, \mathcal{V}_2$ for two different merchants sharing identical counter values $c$.

\subsection{Pre-emption of Delayed-Settlement Insolvency}
\textbf{Theorem 3 (Insolvency Elimination via POE):} Let $B_{\text{offline}}$ be the pre-allocated escrow balance and $\sum a_i$ be the cumulative offline expenditure. Because the TEE locally enforces $\sum a_i \le B_{\text{offline}}$, the clearinghouse $\mathcal{S}$ is guaranteed to hold sufficient escrow reserves to settle all validly signed vouchers upon reconnection.

\subsection{Order-Independent Eventual Consistency}
\textbf{Theorem 4 (Commutative Batch Ingestion):} Let $\mathcal{B}_A$ and $\mathcal{B}_B$ be two disjoint voucher batches containing valid vouchers with non-overlapping counter sets $C_A \cap C_B = \emptyset$. Because set insertion is commutative:
\begin{equation}
(\text{State} \cup C_A) \cup C_B = (\text{State} \cup C_B) \cup C_A
\end{equation}
The final settled state of the clearinghouse ledger is identical regardless of which merchant reconnects first.

---

\section{Experimental Evaluation}

\subsection{Evaluation Testbed \& Cryptographic Latency}
The central switch and client cryptographic primitives were evaluated on an Intel Core i5-8250U workstation @ 1.60\,GHz (4 physical cores, 8 threads) with 16\,GB RAM running Node.js v24 Web Crypto API backed by the OpenSSL cryptographic engine. Edge POS performance was profiled using cycle-accurate mobile silicon benchmarks representative of quad-core ARM Cortex-A53 processors operating at 1.5\,GHz (the dominant architecture in Nigerian agency terminals).

\begin{table}[htbp]
\caption{Cryptographic Execution Latency (ECDSA P-256 / SHA-256)}
\begin{center}
\begin{tabular}{lcccc}
\toprule
\textbf{Cryptographic Operation} & \textbf{Mean (ms)} & \textbf{p50 (ms)} & \textbf{p95 (ms)} & \textbf{p99 (ms)} \\
\midrule
Key Generation & 0.213 & 0.201 & 0.285 & 0.380 \\
Voucher Signing & 0.111 & 0.106 & 0.149 & 0.223 \\
Voucher Verification & 0.152 & 0.146 & 0.186 & 0.243 \\
\bottomrule
\end{tabular}
\label{tab:crypto_latencies}
\end{center}
\end{table}

As shown in Table~\ref{tab:crypto_latencies}, digital voucher verification requires only \textbf{0.152\,ms} (p95: 0.186\,ms) on switch silicon, yielding a single-core throughput exceeding \textbf{6,500 operations/second}. On low-power edge mobile processors (ARM Cortex-A53), ECDSA P-256 verification executes in \textbf{16\,ms to 24\,ms}, providing imperceptible latency for in-person retail checkouts ($<$50\,ms).

\subsection{Payload Serialization Efficiency}
Table~\ref{tab:payload_compaction} highlights the efficiency gains of ResiliPay-Compact over naive JSON serialization.

\begin{table}[htbp]
\caption{Payload Serialization \& Transmission Footprint}
\begin{center}
\resizebox{\columnwidth}{!}{
\begin{tabular}{lccc}
\toprule
\textbf{Serialization Strategy} & \textbf{Raw Bytes} & \textbf{Encoded Length} & \textbf{QR / SMS Target} \\
\midrule
Naive JSON + JWK Base64 & 604 B & 808 chars & QR Version 14 \\
\textbf{ResiliPay-Compact (Binary)} & \textbf{125 B} & \textbf{167 chars} & \textbf{QR Version 6 / 8} \\
\textbf{ResiliPay Single-SMS} & \textbf{86 B} & \textbf{116 chars} & \textbf{1x GSM SMS ($<$160c)} \\
\bottomrule
\end{tabular}
}
\label{tab:payload_compaction}
\end{center}
\end{table}

Condensing JSON into structured binary achieves a \textbf{72.4\% reduction in payload footprint} (from 604 bytes down to 125 bytes). This allows the optical matrix to scale down from Version 14 (73$\times$73 modules) to \textbf{Version 6 or 8} ($41\times41$ to $49\times49$ modules), drastically accelerating camera scan acquisition and improving tolerance to lens blur and screen scratches.

Furthermore, our 86-byte SMS payload (116 Base64 characters) fits comfortably within a single 160-character GSM message, eliminating multi-part SMS concatenation costs.

\subsection{Out-of-Order Burst Reconciliation \& Adversarial Detection}
To validate protocol robustness during chaotic post-blackout reconnection, we simulated 50 offline merchants submitting 1,000 valid queued vouchers in completely randomized, out-of-order batches alongside \textbf{50 injected adversarial double-spend and replay attacks}.

\begin{itemize}
    \item \textbf{Total Vouchers Ingested:} 1,050
    \item \textbf{Valid Out-of-Order Vouchers Settled:} 1,000 (100\% settlement precision)
    \item \textbf{Adversarial Attacks Rejected:} 50 (100\% true-positive rejection)
    \item \textbf{False-Positive Drops:} 0 (0\% false rejection of valid nonces)
    \item \textbf{Total Clearing Ingestion Time:} 176.08\,ms
    \item \textbf{Sustained Switch Throughput:} \textbf{5,679 settled transactions/second}
\end{itemize}

The sustained end-to-end clearing throughput of 5,679 settled tx/sec remains close to the raw cryptographic verification ceiling ($\sim$6,578 ops/sec) because full batch ingestion only incurs minimal $O(1)$ bitmap and escrow-ledger update overhead, while duplicate replay attacks are pre-rejected in $O(1)$ time via bitmap lookup before entering the elliptic curve verification pipeline.

---

\section{Pragmatic Hardware Deployment in Emerging Markets}

To bridge the gap between academic security models and real-world smartphone hardware in Sub-Saharan Africa, ResiliPay introduces a \textbf{Three-Tier Hardware Security Model}:

\subsection{Tier 1: High Assurance (Dedicated Secure Elements)}
\begin{itemize}
    \item \textbf{Hardware:} Google StrongBox (Titan M2) / Apple Secure Enclave.
    \item \textbf{Target Market:} High-end and flagship smartphones.
    \item \textbf{Security Guarantee:} Discrete silicon tamper resistance, physical counter isolation, and high-value offline transactions ($>$ $\mathbb{N}$50,000).
\end{itemize}

\subsection{Tier 2: Standard Mass-Market Assurance (ARM TrustZone + RPMB)}
\begin{itemize}
    \item \textbf{Hardware:} ARM TrustZone Keymaster backed by a Replay Protected Memory Block (RPMB) partition on eMMC/UFS storage.
    \item \textbf{Target Market:} Budget smartphones (\$40--\$90 tier: MediaTek Helio G-series, Unisoc SC9863A/T606 powering Tecno, Infinix, and itel devices dominating 70\%+ of the Nigerian market).
    \item \textbf{Security Guarantee:} Cryptographic rollback resistance enforced by the RPMB shared key written during factory provisioning, supporting medium-value retail transactions (up to $\mathbb{N}$20,000).
    \item \textbf{OEM Integration Path:} Because standard Android user-space applications cannot deploy custom TrustZone Trustlets directly on OEM-locked consumer phones, Tier 2 integration is realized via either: (a) a standardized central-bank/NIBSS Trustlet pre-provisioned in OEM factory firmware through partnerships with major African device distributors (such as Transsion Holdings), or (b) a SIM/eUICC JavaCard applet distributed over-the-air by mobile network operators.
\end{itemize}

\subsection{Tier 3: Software Fallback with Bank Micro-Float}
\begin{itemize}
    \item \textbf{Hardware:} Legacy Android devices lacking RPMB access.
    \item \textbf{Security Guarantee:} The issuing bank underwrites an automated micro-float (e.g., up to $\mathbb{N}$2,000 daily spend) tied to verified national identity (BVN/NIN), contractually absorbing credit default risk to ensure universal financial inclusion across unbanked populations.
\end{itemize}

---

\section{Related Work}
Classic offline payment systems, originating with Chaumian blind signatures~\cite{chaum1983blind}, focused on anonymity and post-facto double-spender identification. However, unbacked post-facto discovery is unacceptable in cash-dominant emerging economies where debt recovery from anonymous consumers is legally and practically impossible.

Modern Central Bank Digital Currency (CBDC) research from the Bank for International Settlements (BIS) and MIT Project Hamilton has explored high-throughput centralized architectures~\cite{bis2023polaris, chu2022offline, lovejoy2023hamilton}. Visa's Offline Payment System (OPS) proposed protocol mechanisms relying on pre-funded allowances and secure hardware~\cite{christodorescu2020towards}. ResiliPay extends this literature by solving out-of-order multi-merchant batch settlement, optimizing dual-channel transmission (QR + SMS) for budget consumer hardware, and providing empirical validation tailored to the infrastructural realities of Sub-Saharan Africa.

---

\section{Conclusion}
This paper presented \textbf{ResiliPay}, a fault-tolerant, decentralized offline payment protocol engineered to solve digital payment fragility across intermittently connected economies. By coupling ECDSA P-256 signatures with hardware-enforced monotonic counters and Pre-Allocated Offline Escrow within commodity mobile TEEs, ResiliPay eliminates both offline double-spending and delayed-settlement insolvency without requiring distributed consensus. Our compact binary serialization achieves a 72.4\% reduction in QR payload size and enables wide-area cellular fallback within a single 160-character GSM SMS segment. With an order-independent reconciliation clearing engine processing over 5,600 transactions/second and a pragmatic three-tier hardware model, ResiliPay provides an actionable, accessible foundation for sovereign financial resilience in the developing world.

\bibliographystyle{IEEEtran}
\bibliography{references}

\end{document}
